% !PRINTED 610
\weberpage{610}
%<<CONT
Gleichung durch Radicale zu gestatten. Wir leiten also, zunächst unter der
Voraussetzung, dass der Grad $n$ der Gleichung eine Primzahl sei, ein anderes von
\gesperrt{Galois} zuerst aufgestelltes Kriterium her.

Es sei jetzt $f(x) = 0$ eine irreducible Gleichung vom Grade $n$ und $n$ eine Primzahl
grösser als 2. Soll $f(x)$ algebraisch lösbar sein, so muss nach \S.\,176 ihre Gruppe
$P$ metacyklisch sein, d. h. es muss eine Kette von Gruppen
\weq{1}{P,\ P_1,\ P_2 \ldots P_{\mu-1},\ 1}
geben, deren jede ein Normaltheiler der nächst vorangehenden mit Primzahlindex ist.
Durch successive Adjunction von Wurzeln cyklischer Gleichungen von Primzahlgrad wird
die Gruppe der Gleichung von $P$ auf $P_1, P_2 \ldots P_{\mu-1}, 1$ reducirt.

Die Function $f(x)$ selbst kann, wie schon im \S.\,176 gezeigt ist, nicht reducirt
werden, ehe die letzte Adjunction gemacht ist, worauf sie in $n$ lineare Factoren
zerfällt. Der Grad der letzten cyklischen Gleichung und mithin der Grad der vorletzten
Gruppe $P_{\mu-1}$ muss also nach \S.\,158, 3. gleich $n$ sein.

Der Grad eines Elementes einer Gruppe ist immer ein Theiler vom Grade der Gruppe und
daher kann $P_{\mu-1}$ ausser der identischen nur Permutationen von der Ordnung $n$
enthalten. Sie ist also mit der Periode $1, \pi, \pi^2 \ldots \pi^{n-1}$ eines ihrer
Elemente identisch. Daraus folgt nach \S.\,153, dass $\pi$ eine cyklische Permutation
von sämmtlichen $n$ Ziffern sein muss, und wir können die Bezeichnung der Wurzeln von
$f(x)$ so wählen, dass
\[ \pi = (0,\ 1,\ 2 \ldots n-1) \]
wird. Bezeichnen wir die Wurzeln mit $x_z$, und setzen fest, dass $x_z = x_{z'}$ sein
soll, wenn $z \equiv z' \, (\mathrm{mod}\ n)$ ist, so geht durch $\pi$ jedes $z$ in
$z+1$ über, durch $\pi^2$ in $z+2$ u. s. f., und wir können also sagen, dass die
Gruppe $P_{\mu-1}$ aus den Substitutionen für $z$
\weq{2}{(z,\ z+b), \qquad b = 0,\ 1,\ 2 \ldots n-1}
besteht. Jede irreducible metacyklische Gleichung von Primzahlgrad kann also durch
Adjunction von Radicalen in eine cyklische Gleichung verwandelt werden.

Die Substitutionen $(z, z+b)$ bilden einen speciellen Fall der allgemeineren linearen
Substitution
\weq{3}{(z,\ az+b),}
worin $a$ und $b$ feste Zahlen sind, die auch nach dem Modul $n$ reducirt werden
können, wo aber für $a$ der Werth 0 natürlich
%CONT>>
